\sscp{\tsc{Sumba-Havu} as a subgroup}{03-03-03}
\citet{Blust2008} proposes a \tsc{Sumba-Havu} node based on four
lexically specific phonological irregularities,
six instances of morphological fossilization,
and 97 shared lexical items.
Blust purposely limits his investigation to Kambera
and Havu as the best-described languages of this node.
Examination of other putative members of \tsc{Sumba-Havu} shows
that the proposed phonological irregularities\footnote{
		\citet[90]{Blust2008}
		proposes the following lexically specific phonological irregularities:
		(i.) Irregular *ə > **o in *qaləjaw > putative Proto-Sumba-Havu (PSH) **loɗo.
		This cannot account for Mangili,
		Anakalang, Mamboru ‹làdo› [lɐɗo] or Laboya \it{laɗo}.
		(ii.) *p > **ʔ in *dapuR > PSH **raʔu ‘hearth’.
		This cannot account for Wanukaka \it{rapu} ‘hearth’.
		(iii.) Consonant metathesis in *qulun-an > PSH **nulaŋ ‘wooden headrest, pillow’.
		This cannot account for Weyewa \it{luna}.
		(iv.) *añəm > PSH **uñaŋ.
		This cannot account for Anakalang \it{wanaŋ-u},
		Mamboru, Weyewa \it{wanan-a}, Kodi \it{ɣanaŋ-o},
		or Loli \it{nane}.} 
and all but one of the morphological irregularities\footnote{
		\citet[91]{Blust2008}
		proposes the following morphological irregularities: (i.) PMP
		*kita ‘see’ > Proto-Sumba-Havu **ŋita,
		for which Blust concedes both \it{ita}
		and \it{ŋita} occur for Kambera.
		Putative **ŋita is further not known in Dhao or other languages of Sumba.
		(ii.) *qunap > PSH **napi ‘fish scale’ has probable cognates
		in Rote, such as Termanu \it{hani-k} ‘shell’ (with consonant metathesis from intermediate **pani).
		(iii.) PMP *di atas > PSH **diatas ‘above,
		on top’ has cognates in Flores (e.g.
		Central Flores has **deta ‘above/mountain-wards’ -- see \trf{03-37}
		where Proto-Nuclear Central Flores has **ʤeta).
		(iv.) *qəti > PSH **mənti ‘used up’ would be valid.
		(v.) *laRiw > PSH **pa-lai has cognates in \tsc{Flores-Lembata} (e.g.
		see reflexes \trf{03-46} \prf{tab:03-46})
		and Blust concedes that forms without initial *pa- occur in both Kambera and Havu.
		(vi.) *maŋ-kaən > PSH **ŋaŋ ‘eat’.
		As Blust discusses, this form can be derived from PMP with loss
		of a morpheme boundary.
		It also does not account for Dhao \it{-aʔa} ‘eat’.}
		do not account
for a wider range of data.
This would leave a single morphological innovation (*qəti > **mənti
‘used up, finished’) and the putative lexical innovations as evidence for \tsc{Sumba-Havu}.
Of the 97 putative lexical innovations,
Blust notes that eighteen have problematic vowel correspondences.
We further note that eight have cognates elsewhere,\footnote{
		We have identified external cognates %in Timor
		for the following
		putative Proto-Sumba-Havu (PSH) reconstructions in \citet[109--111]{Blust2008}:
		PSH **kərasa ‘ribs’ (cf.
		Proto-Rote-Meto [PRM] **kara ‘chest’ from earlier **karas (consonant
		metathesis) and/or **ŋgarasa ‘part of the back’),
		PSH **ləmisi ‘pit, seed of fruit’ (cf.
		PRM *muse ‘seed, pip’ which is a better match for the penultimate
		vowel of Havu \it{lamuhi},
		as well as Tetun \it{musan} ‘seed,
		pip’), PSH **mawo ‘shadow’ (cf.
		PRM **mafo ‘cool, shadow’ and other cognates including
		Bima \it{mawo} ‘cool down,
		no longer hot, shady’ [Ismail et al.
		1985:87]), PSH **muku ‘banana’ (many cognates throughout the
		Lesser Sundas, see \srf{02-03-01-07}, \srf{15-07}, and appendix
		\ref{ch:BanTer}), PSH **ŋəlu ‘wind’ (many cognates widely distributed,
		see \srf{15-04-02-02}), PSH **samba ‘bucket’ (cf.
		PRM **sambat ‘bucket’), PSH **tənaŋga ‘anchor’ (cf.
		Termanu \it{naka}, Rikou \it{naʔa} ‘anchor’), and PSH **tuku ‘throw’ (cf.
		PRM **tuku ‘throw’).} 
and six are inherited from PMP.\footnote{
		Putative Proto-Sumba-Havu
		(PSH) reconstructions in \citet[109--111]{Blust2008} which are
		derived through regular processes from PMP are: PSH **ae ‘much’
		< *Raya ‘big’, PSH **la ‘towards’ reduction of PMP *lakaw (cf.
		Weyewa \it{laʔo}, \it{laʔa} ‘go’),
		PSH **liu ‘behind’ < *likud ‘back’ (putative **liu also does
		not account for Dhao \it{liʔu} ‘back,
		behind, outside’), PSH **mami ‘ripe’ < *ma-həmis ‘sweet’ (noted
		as a possibility by Blust),
		*məndua ‘half’ < *ma- + *duha ‘two’,
		and PSH **riŋu ‘cold’ for which Havu \it{meriŋi} ‘cold’
		and Dhao \it{riŋi} ‘cold’ are regular from *ma-diŋdiŋ.}

\largerpage
While we are sceptical of subgroups based on lexical evidence
(see \srf{15-04}),
it is worth investigating whether there is
additional evidence in favour of \tsc{Sumba-Havu}.
This subgroup is supported on the basis of at least
two shared sound changes: *R > **y > {\0} in a subset of words
and *j > **ɗ, **r (both identified by \citealp[93--94]{Blust2008}).
%\footnote{
		%\citet{Blust2008} identifies both *R > {\0}
		%and *d/*j > \it{r} for Kambera
		%and Havu, but appears to treat them as cases of parallel development (cf.
		%\citealp[93--94]{Blust2008}).}
Within \tsc{Bima-Lembata}
*R > **y > {\0} is mostly unique,
though \tsc{East Sumbawa}, Komodo (\srf{03-01-01}),
and some members of \tsc{West Lamaholot} (\srf{03-02-04}) have similar *R > {\0}.

\tsc{Havu-Dhao} have *R > {\0} in nearly all words,
while languages of Sumba have a split with *R > {\0}, \it{r}.
Nonetheless, there are eleven words with cognates in both Sumba
and \tsc{Havu-Dhao} which share *R > {\0}.
These words are given in \trf{03-51}.
Words which provide evidence of earlier **y are shaded.

\begin{table}[p]\centering\tcp{\tsc{Sumba-Havu} non-final *R > {\0}}{03-51}
%\begin{adjustbox}{width=\textwidth}
	\begin{threeparttable}\st{2.75pt}
	\begin{tabular}{llllllll}\lsptoprule
*gloss	&	root	&	house	&brave	&	heavy	&	blood	&	red\su{†}	&	thorn\su{‡}\\
PMP	&	*Ramut	&	*Rumaq	&*baRani	&	*bəRəqat	&	*daRaq	&	*ma-iRaq	&	*duRi\\	\midrule
Kodi	&	\hp{*R}amu	&	\hp{*R}umːa	&	{\hr}mbani	&	{\hr}mboto	&	\cg{(rutːa)}	&		&	{\hr}riyo\su{Δ}\\
Weyewa	&		&	\hp{*R}umːa	&	{\hr}mbani	&	{\hr}mboto	&	{\hr}raʔa	&	{\hr}meta	&	{\hr}ruwi\\
Loli	&		&	\hp{*R}umːa		&	&	{\hr}buat-o	&	{\hr}raʔa	&	{\hr}miet-a	{\cgr}&	{\hr}rii{\cgr}\\
Laboya	&		&	\hp{*R}uma		&	&	{\hr}biʔta	{\cgr}&	{\hr}raa	&		&	{\hr}rii{\cgr}\\
Wanukaka	&		&	\hp{*R}uma	&	{\hr}bani	&	{\hr}buat-u	&	{\hr}raa, raʔa	&	{\hr}miaŋ	{\cgr}&	{\hr}rii{\cgr}\\
Mamboru	&	\hp{*R}amu	&	\hp{*R}uma	&	{\hr}mbani	&	{\hr}mbuat-u	&	{\hr}ree	{\cgr}&		&	{\hr}rii{\cgr}\\
Anakalang	&	\hp{*R}amu	&	\hp{*R}uma	&	{\hr}bani	&	{\hr}buat-u	&	{\hr}ree	{\cgr}&		&	{\hr}rii{\cgr}\\
Lewa	&	\hp{*R}amu	&	\hp{*R}uma	&	{\hr}mbanu	&	{\hr}mbuat-u	&	{\hr}ria	{\cgr}&	{\hr}mia	{\cgr}&	{\hr}rii{\cgr}\\
Kambera	&	\hp{*R}amu	&	\hp{*R}uma	&	{\hr}mbeni	&	{\hr}mbotu	&	{\hr}ria	{\cgr}&	{\hr}mii	{\cgr}&	{\hr}rii{\cgr}\\
Mangili	&	\hp{*R}amu	&	\hp{*R}uma	&	{\hr}mbani	&	{\hr}mbotu	&	{\hr}ria	{\cgr}&	{\hr}mii	{\cgr}&	{\hr}rii{\cgr}\\
Havu	&	\hp{*R}amo	&	\hp{*R}əmu	&	{\hr}ɓani	&	\cg{(meʤəni)}	&	{\hr}raa	&	{\hr}mea	&	{\hr}rui\\
Dhao	&	\hp{*R}amo	&	\hp{*R}əmu	&	{\hr}bani	&	{\hr}bia	{\cgr}&	{\hr}raa	&	{\hr}mea	&	{\hr}rui\\
\midrule\midrule
\end{tabular}
	\begin{tabular}{lllllll}
*gloss	&	run	&	scratch	&	bite	&	octopus	&	giant taro	&	rotten\\
PMP	&	*laRiw	&	*kaRaw	&	*kaRat-i	&	*kuRita	&	*biRaq	&	*baRiw\\	\midrule
Kodi	&	{\hr}palayo\su{Δ}	&	{\hr}kayo	&	{\hr}kati	&	{\hr}kawicːa	&	{\hr}wiyo	&	\cg{(mbeŋːo)}\\
Weyewa	&	{\hr}malːe	&	{\hr}kau	&	{\hr}kati	&	{\hr}kawitːa	&	{\hr}wiʔa	&	\cg{(mbeŋːo)}\\
Loli	&	{\hr}malːe	&	{\hr}kaʔu	&	{\hr}pakati	&		&		&	\\
Laboya	&	{\hr}malai	&	{\hr}kau	&	{\hr}kaːti	&		&		&	\\
Wanukaka	&	{\hr}palai	&	\cg{(gowa)}	&	{\hr}kahi	&		&		&	\\
Mamboru	&	{\hr}palai	&	{\hr}kau	&	{\hr}kati	&	{\hr}kawita	&	{\hr}wia	&	{\hr}mbai\\
Anakalang	&	{\hr}palai	&	{\hr}kau	&		&	{\hr}mawita	&	{\hr}wia	&	{\hr}bai\\
Lewa	&	{\hr}palee	&	{\hr}kau	&		&	{\hr}wita, kawita	&	{\hr}wii	&	{\hr}mbai\\
Kambera	&	{\hr}palai	&	{\hr}kau	&	{\hr}kati	&	{\hr}wita, kawita	&	{\hr}wii	&	{\hr}mbai\\
Mangili	&	{\hr}palai	&	{\hr}kau	&	{\hr}kati	&	{\hr}wita, kawita	&	{\hr}wii	&	{\hr}mbia\\
Havu	&	{\hr}perai	&	{\hr}kao, kau	&		&	{\hr}ədi	&		&	\\
Dhao	&	{\hr}rai	&	{\hr}kao	&	{\hr}kaɖ͡ʐi	&	\cg{(kapaʄu)}	&		&	\\
		\lspbottomrule
	\end{tabular}
		\begin{tablenotes}\tnsize
			\item[†]	{Loli \it{miet-a} is glossed ‘solid red’ (Indonesian ‘bewarna
								merah pekat’) and Wanukaka \it{miaŋ} as ‘reddish’.
								Havu and Dhao \it{mea} are ‘reddish-brown’ Other reflexes of
								*ma-iRaq are from \citet[287]{Onvlee1984} who defines Kambera
								\it{mii} as ‘red (of lips)’.}
			\item[‡]	{Anakalang,
								Lewa, Kambera, and Mangili reflexes of *duRi mean ‘thorn, bone’.
								The only known meaning for the other reflexes here is ‘bone’,
								though Kodi has cognate \it{kariyo} ‘thorn’.}
			\item[Δ]	{Kodi \it{riyo} ‘bone’
								and \it{palayo} ‘run’ have insertion of a glide in historic vowel sequences.
								Two other examples include *buah > \it{wuyo} ‘fruit’,
								and *tau > \it{toyo} ‘person’.}
			%\item[‖]	{Kambera also has the word \it{lai} which occurs in the phrase
								%\it{ɲʤara lai} ‘racing horse’.}
		\end{tablenotes}
	\end{threeparttable}
%\end{adjustbox}
\end{table}

The shaded words in \trf{03-51} show palatalisation or raising
of vowels adjacent to *R.
This provides indirect evidence that *R went through an intermediate
glide stage in at least some words.
This change -- loss with palatalisation of adjacent vowels ---
is consistent with the changes affecting PMP *y in Sumba and \tsc{Havu-Dhao}.
Word medially PMP *y > \it{y,} {\0}[+fr] in Sumba
and *y > {\0} in \tsc{Havu-Dhao}.
Two examples of *y > {\0}[+fr] in Sumba are *daya ‘upstream’
>  Kambera,
Mangili, Lewa \it{ɗia}, Anakalang \it{ɗee} (for which the changes
affecting the vowels are identical to those for *daRaq ‘blood’)
and *layaR ‘sail’ > Kambera,
Mangili, Lewa, Anakalang \it{liru} ‘sail’, Mamboru,
Wanukaka \it{leru} ‘sail’, Weyewa,
Kodi \it{lere}, and Laboya \it{leːro} (all with a historically epenthetic final vowel).
\citet[69]{Blust2008} also proposes that *R > {\0} in Havu went
through intermediate **y due to otherwise unexpected *ə > \it{i}
in reflexes of *baqəRu ‘new’ > Havu \it{viu,} Dhao \it{hiu};
*dəŋəR ‘hear’ > Havu \it{rəŋe {\tl} rəŋi} ‘hear', Dhao \it{tadəŋi} ‘hear’;
and *bəRay ‘give’ > Havu \it{vie}, Dhao \it{hia}.

Some words with medial *R have a medial glottal stop in Weyewa,
Loli, Laboya, and Wanukaka, which could be a reflex of *R.
However, these languages have a more general process of glottal stop insertion.
Examples include: *dahun > Weyewa,
Loli \it{roʔo} ‘leaf’ (Kambera \it{rau,
ruu}), *bayu > Weyewa, Loli \it{baʔi} ‘pound’ (Kambera \it{ɓai}),
*kahu > Weyewa \it{woʔu} ‘\tsc{2sg}’ (Loli \it{wu,} Wanukaka
\it{ou)}, and *daqan > Weyewa \it{karaʔa} ‘branch’.
Thus, these medial glottal stops are not due to *R > \it{ʔ},
but rather due to *R > {\0} followed by {\0} > \it{ʔ} /V{\gap}V.

\clearpage
While \tsc{Sumba-Havu} share *R > {\0} in some words,
there are also words in which *R = \it{r} in the languages of Sumba.
Sometimes there is a split of *R > \it{r,} {\0} between different
languages -- or even within a single language ---
and sometimes *R = \it{r} in all known reflexes.
Examples are given in \trf{03-53}.
Havu and Dhao have *R > {\0} in all such words when cognates are available.

\begin{table}\tcp{\tsc{Sumba-Havu} retentions of *R}{03-53}
%\begin{adjustbox}{width=\textwidth}
	\begin{threeparttable}\st{3.5pt}
	\begin{tabular}{lll@{ }ll@{ }lll@{ }l}\lsptoprule
*gloss	&	new	&	\mc{2}{l}{sea hibiscus}			&	\mc{2}{l}{ray, skate}			&	west (wind)	&	\mc{2}{l}{artery, vein}		\\
PMP	&	*baqə\tbr{R}u	&	\mc{2}{l}{*ba\tbr{R}u}			&	\mc{2}{l}{*pa\tbr{R}ih}			&	*haba\tbr{R}at	&	\mc{2}{l}{*u\tbr{R}at\su{‡}}		\\	\midrule
Kodi	&	{\hr}mba\tbr{r}u	&	\mc{2}{l}{{\hr}kawaŋ-o}			&	\mc{2}{l}{{\hr}paya}			&	\hp{*ha}wa\tbr{r}at-o	&	\mc{2}{l}{{\hr}u\tbr{r}at-o}		\\
Weyewa	&	{\hr}mba\tbr{r}u	&	\mc{2}{l}{{\hr}kawaŋ-o}			&	\mc{2}{l}{{\hr}pai, pawe}			&	\hp{*ha}wa\tbr{r}at-a	&	{\hr}u\tbr{r}a,	&	u\tbr{r}at-a\\
Loli	&	{\hr}boʔu	&		&		&		&		&		&	\mc{2}{l}{{\hr}u\tbr{r}at-a}		\\
Laboya	&	{\hr}beː\tbr{r}i	&		&		&		&		&		&	\mc{2}{l}{\cg{(kaluta)}}		\\
Wanukaka	&	{\hr}ɓua	&	\mc{2}{l}{{\hr}kawauŋ-u}			&	{\hr}pai	&		&	\hp{*ha}wa\tbr{r}at-u\su{†}	&	{\hr}u\tbr{r}a,	&	u\tbr{r}at-u\\
Mamboru	&	{\hr}mba\tbr{r}u	&	\mc{2}{l}{{\hr}kawau}			&	{\hr}pai	&		&	\hp{*ha}wa\tbr{r}at-a	&	{\hr}u\tbr{r}a,	&	u\tbr{r}at-u\\
Anakalang	&	{\hr}bua	&	\mc{2}{l}{{\hr}kawauŋ-u}			&	{\hr}pa\tbr{r}i	&		&		&	{\hr}u\tbr{r}a,	&	u\tbr{r}at-u\\
Lewa	&	{\hr}mba\tbr{r}u	&	{\hr}wa\tbr{r}u,	&	kaɓa\tbr{r}u	&	{\hr}pa\tbr{r}i	&		&	\hp{*ha}wa\tbr{r}at-u	&	{\hr}u\tbr{r}a,	&	u\tbr{r}at-u\\
Kambera	&	{\hr}mba\tbr{r}u	&	{\hr}wa\tbr{r}u,	&	kaɓa\tbr{r}u	&	{\hr}pa\tbr{r}i,	&	pai	&	\hp{*ha}wa\tbr{r}at-u	&	{\hr}u\tbr{r}a,	&	u\tbr{r}at-u\\
Mangili	&	{\hr}mba\tbr{r}u	&	{\hr}wa\tbr{r}u,	&	kaɓa\tbr{r}u	&	{\hr}pa\tbr{r}i,	&	pai	&	\hp{*ha}wa\tbr{r}at-u	&	{\hr}u\tbr{r}a,	&	u\tbr{r}at-u\\
Havu	&	{\hr}viu	&	\mc{2}{l}{{\hr}vau}			&	{\hr}ai	&		&	\hp{*ha}vaa	&	\mc{2}{l}{\cg{(lua)}}		\\
Dhao	&	{\hr}hiu	&	\mc{2}{l}{{\hr}hau}			&	{\hr}ai	&		&	\hp{*ha}haa	&	\mc{2}{l}{\cg{(kalua)}}		\\
		\lspbottomrule
	\end{tabular}
		\begin{tablenotes}\tnsize
			\item[†]	{Wanukaka \it{warat-u} is glossed ‘gale’.}
			\item[‡]	{Kambera \it{ura} = ‘(hand)lines,
								decision, fate, power’
								%(Dutch `[hand]lijnen, beschikking, lot, kracht'),
								and \it{urat-u} = ‘line, inspect lines, predict the future’.
								Wanukaka \it{ura} is given as ‘fingerprints; lines or marks that
								show a person’s identity’
								and \it{urat-u} as ‘divination; with spear pressed on at main
								house pillar, gong beaten’.
								Loli \it{urat-a} is glossed ‘pray’.}
		\end{tablenotes}
	\end{threeparttable}
%\end{adjustbox}
\end{table}

\clearpage
An additional change shared between Sumba
and \tsc{Havu-Dhao} is of *j > \it{ɗ} /*ə{\gap}
and *j > \it{r} elsewhere.
Examples are given in \trf{03-54}.
One exception to *j > \it{r} is *huaji ‘same sex younger sibling’
which has *j > \it{l} in Weyewa and Loli \it{alli}.
This word has *j > \it{r} in other known reflexes in \tsc{Sumba-Havu}:
Kodi, Anakalang, Mamboru, Mangili, Havu,
Dhao \it{ari}, Lewa \it{aru},
and Kambera \it{eri}.

\begin{table}\tcp{\tsc{Sumba-Havu} *j > **ɗ, \it{r}}{03-54}
\begin{adjustbox}{width=\textwidth}
	\begin{threeparttable}\st{2pt}
	\begin{tabular}{lllllll@{}l}\lsptoprule
*gloss	&	day, sun		&	gall, bile			&	name		&	rice		&	how many			&	foam		&	nose		\\
PMP	&	*qalə\tbr{j}aw	&	*qapə\tbr{j}u\su{†}			&	*ŋa\tbr{j}an		&	*pa\tbr{j}ay		&	*pi\tbr{j}a			&	*bu\tbr{j}əq		&	*i\tbr{j}uŋ/*u\tbr{j}uŋ	\\	\midrule
Kodi	&	{\hqa}lo\tbr{ɗ}ːo	&	{\hqa}po\tbr{ɗː}u		&	{\hr}ŋa\tbr{r}a	&	{\hr}pa\tbr{r}i	&	{\hr}pi\tbr{r}a\su{‡}		&	‹wu\tbr{r}a›		&	{\hr}i\tbr{r}u	\\
Weyewa	&	{\hqa}lo\tbr{ɗ}ːo	&	{\hqa}po\tbr{ɗ}ːu		&	{\hr}ŋa\tbr{r}a	&	{\hr}pa\tbr{r}e	&	{\hr}pi\tbr{r}ːa		&	‹wu\tbr{r}a›		&	{\hr}ko\tbr{r}ːu\su{ǁ}	\\
Loli	&	\hp{qa}‹lo\tbr{dd}o›	&	\hp{qa}‹po\tbr{dd}u›		&	{\hr}ŋa\tbr{r}a	&	{\hr}pa\tbr{r}e	&	‹pi\tbr{r}a, pi\tbr{rr}a›\su{Δ}	&	‹wu\tbr{rr}a›		&	‹ko\tbr{rr}u›\su{ǁ}	\\
Laboya	&	{\hqa}la\tbr{ɗ}o	&	{\hqa}pa\tbr{ɗ}u		&	{\hr}ŋaː\tbr{r}a	&	{\hr}paː\tbr{r}e	&	{\hr}pi\tbr{r}a		&	kawuː\tbr{r}akana	&	{\hr}ka\tbr{r}u\su{ǁ}	\\
Wanukaka	&	{\hqa}la\tbr{d}u	&	{\hqa}pi\tbr{d}u		&	{\hr}ŋa\tbr{r}a	&	{\hr}pa\tbr{r}i	&	{\hr}pi\tbr{r}a		&	‹wu\tbr{rr}a›		&	{\hr}i\tbr{r}u	\\
Mamboru	&	{\hqa}lɐ\tbr{ɗ}u	&	{\hr}kapɐ\tbr{ɗ}u		&	{\hr}ŋa\tbr{r}a	&	{\hr}pa\tbr{r}i	&	{\hr}pi\tbr{r}a		&	{\hr}wu\tbr{r}a	&	{\hr}u\tbr{r}u	\\
Anakalang	&	\hp{qa}‹là\tbr{d}u›	&	{\hr}kapi\tbr{d}u		&	{\hr}ŋa\tbr{r}a	&	{\hr}pa\tbr{r}i	&	{\hr}pi\tbr{r}a		&	{\hr}wu\tbr{r}a	&	{\hr}ki\tbr{r}u\su{ǁ}	\\
Lewa	&			&	{\hr}kapi\tbr{ɗ}u		&	{\hr}ŋa\tbr{r}a	&	{\hr}pa\tbr{r}i	&	{\hr}pi\tbr{r}a		&	{\hr}wu\tbr{r}a	&	{\hr}i\tbr{r}u	\\
Kambera	&	{\hqa}lo\tbr{ɗ}u	&	{\hr}kapi\tbr{ɗ}u		&	{\hr}ŋa\tbr{r}a	&	{\hr}pa\tbr{r}i	&	{\hr}pi\tbr{r}a		&	{\hr}wu\tbr{r}a	&	{\hr}u\tbr{r}u	\\
Mangili	&	{\hqa}lɐ\tbr{ɗ}u	&	{\hr}kapɐ\tbr{ɗ}u		&	{\hr}ŋa\tbr{r}a	&	{\hr}pa\tbr{r}i	&	{\hr}pi\tbr{r}a		&	{\hr}wu\tbr{r}a	&	{\hr}u\tbr{r}u	\\
Havu	&	{\hqa}lo\tbr{ɗ}o	&	{\hr}(p)ə\tbr{ɗ}u	&	{\hr}ŋa\tbr{r}a	&	\hp{*p}a\tbr{r}e	&	{\hr}pə\tbr{r}i		&	{\hr}vo\tbr{r}o	&	\cg{(hevəŋa)}		\\	
%	&		&	\hp{*qap}ə\tbr{ɗ}u	&		&		&			&		&			\\
Dhao	&	{\hqa}lo\tbr{ɗ}o	&	{\hr}papə\tbr{d}u		&	{\hr}ŋa\tbr{r}a	&	\hp{*p}a\tbr{r}e	&	{\hr}pə\tbr{r}i		&	{\hr}ho\tbr{r}o	&	\cg{(sahəŋa)}		\\
		\lspbottomrule
	\end{tabular}
		\begin{tablenotes}\tnsize
			\item[†]	{Laboya \it{paɗu},
								Wanukaka \it{piɗu}, Loli ‹\it{poddu›},
								Havu \it{pəɗu}, and Dhao \it{papədu} are all ‘bitter’.
								Other words given here (including Havu \it{əɗu}) are glossed
								as ‘gall, bile’ (Dutch \it{gal}).}
			\item[‡]	{\citet[415]{Onvlee1984} gives Kodi ‹piri› ‘how many’.}
			\item[Δ]	{Loli ‹pira› is given as ‘when’
								and ‹pirra› as ‘how many’.
								Whether the distinction between \it{pira}
								and \it{pirra} is purely orthographic or
								reflects a real difference is unknown.}
			\item[{ǁ}]	{Words with initial \it{k} do not
								appear to be straightforward reflexes of *ijuŋ/*ujuŋ.}
		\end{tablenotes}
	\end{threeparttable}
\end{adjustbox}
\end{table}
